% Yêu cầu 5 — Ngân sách: ~3.5 trang
\chapter{Hạn chế và kết luận}\label{chap:hanche}

\section{Hạn chế của phần mềm}\label{sec:hanche}
Ngoài kịch bản chính ở Chương~\ref{chap:demo}, nhóm chạy thêm ba thí nghiệm trên cùng lab để lượng hóa các hạn chế: 15 biến thể payload (\texttt{lab/cases/bypass.txt}), 10 thao tác hợp lệ trên DVWA ở bốn mức paranoia (\texttt{lab/cases/benign.txt}), và đo hiệu năng (\texttt{lab/scripts/bench.sh}). Kết quả lưu trong \texttt{lab/evidence/}.

\subsection{Bỏ sót tấn công và vượt qua bằng biến thể payload}
Bảng~\ref{tab:bypass} cho thấy CRS ở PL1 chống tốt các kỹ thuật né tránh kinh điển cho SQLi, XSS và LFI. Đổi hoa thường, chèn chú thích \texttt{/**/}, thay \texttt{OR} bằng \texttt{||} hay mã hóa URL hai lần đều không qua được, vì trước khi so khớp, ModSecurity giải mã và chuẩn hóa dữ liệu (\texttt{t:urlDecodeUni}, \texttt{t:lowercase}, \ldots), còn libinjection phân tích theo token chứ không theo chuỗi cố định.

\begin{table}[!htb]
\centering
\caption{Kết quả 15 biến thể payload (số bị chặn / số biến thể)}
\label{tab:bypass}
\footnotesize
\begin{tabularx}{\linewidth}{@{}l>{\raggedright\arraybackslash}Xccl@{}}
\toprule
\textbf{Loại} & \textbf{Biến thể} & \textbf{PL1} & \textbf{PL3} & \textbf{Luật chính} \\ \midrule
SQLi & \path{/**/} thay khoảng trắng; \path{||} thay \texttt{OR}; \texttt{LIKE}; mã hóa URL 2 lần & 4/4 & 4/4 & 942100 \\
XSS & \path{<ScRiPt>}; \path{<svg/onload=...>}; \path{<img onerror=...>} & 3/3 & 3/3 & 941100--941390 \\
LFI & \path{/etc/passwd}, \path{/etc/hosts} (không có \texttt{../}); \path{php://filter} & 3/3 & 3/3 & 930120, 933140 \\
CMDi & \path{;cat ../../config/config.inc.php} & 1/1 & 1/1 & 930100--930120 \\
CMDi & \path{;id}, \path{|id}, xuống dòng + \texttt{id}, \path{;pwd} & \textbf{0/4} & 4/4 & 932350 (PL3) \\
\bottomrule
\end{tabularx}
\end{table}

Điểm yếu nằm ở payload ngắn, ít đặc trưng. Bốn biến thể command injection dùng lệnh hai, ba ký tự (\texttt{id}, \texttt{pwd}) lọt qua PL1 hoàn toàn và chạy được trên máy chủ, vì các luật bắt lệnh Unix không kèm đối số chỉ có từ PL3. Đây là hệ quả của triết lý thiết kế CRS: PL1 ưu tiên ít chặn nhầm, chấp nhận bỏ sót. Kẻ tấn công có thể dùng chính các lệnh ngắn này để dò môi trường trước khi viết payload dài hơn.

Ngoài ra, WAF dựa trên chữ ký không hiểu ngữ nghĩa ứng dụng. Nó không bảo vệ được lỗi logic nghiệp vụ, lỗi phân quyền (IDOR) hay CSRF (yêu cầu giả mạo trông hoàn toàn hợp lệ), và không chống dò mật khẩu nếu không kết hợp giới hạn tốc độ. Với lưu lượng HTTPS, WAF chỉ có tác dụng khi được đặt ở điểm kết thúc TLS.

\subsection{Chặn nhầm yêu cầu hợp lệ}
Bảng~\ref{tab:fp} là số thao tác hợp lệ (trong 10 thao tác) bị chặn ở từng mức. Ở PL1 không có chặn nhầm. Từ PL2, một lời nhắn sổ lưu bút có các từ \textit{select \ldots\ from \ldots\ drop} bị luật 942200/942260 coi là SQLi. Ở PL3, luật 932350 chặn mọi yêu cầu có tham số tên \texttt{id} như đã phân tích ở mục~\ref{sec:pl}. Ở PL4, luật 920273 chỉ chấp nhận một tập ký tự rất hẹp nên chặn cả tên \textit{Nguyen Van An} (vì dấu cách), \textit{O'Brien} và một địa chỉ email.

\begin{table}[!htb]
\centering
\caption{Số thao tác hợp lệ bị chặn nhầm (trên 10 thao tác)}
\label{tab:fp}
\footnotesize
\begin{tabularx}{\linewidth}{@{}lcc>{\raggedright\arraybackslash}X@{}}
\toprule
\textbf{Mức} & \textbf{Không loại trừ} & \textbf{Có 1 luật loại trừ} & \textbf{Luật gây chặn nhầm} \\ \midrule
PL1 & 0 & 0 & -- \\
PL2 & 1 & 1 & 942200, 942260 (sổ lưu bút) \\
PL3 & 3 & 1 & 932350 (tham số \texttt{id}); 942200, 942260 \\
PL4 & 6 & 4 & thêm 920273 (ký tự ngoài tập rất hẹp), 942432 \\
\bottomrule
\end{tabularx}
\end{table}

Kết hợp hai bảng cho thấy sự đánh đổi trực tiếp: muốn chặn được \texttt{;id} phải lên PL3, nhưng khi đó 3/10 thao tác hợp lệ bị chặn.

\subsection{Chi phí tinh chỉnh luật}\label{sec:tinhchinh}
Cách khắc phục chặn nhầm là viết luật loại trừ (rule exclusion). Listing~\ref{lst:excl} là luật nhóm viết cho DVWA, nạp sau CRS qua biến \texttt{EXCLUSION\_FILE}. Luật chỉ bỏ \emph{tên} tham số \texttt{id} khỏi danh sách biến mà 932350 kiểm tra; \emph{giá trị} các tham số vẫn được kiểm tra như cũ.

\begin{lstlisting}[caption={Luật loại trừ cho DVWA (\texttt{lab/rules/dvwa-exclusions.conf})},label={lst:excl}]
SecRuleUpdateTargetById 932350 "!ARGS_NAMES:id"
\end{lstlisting}

Với một dòng này, PL3 cho qua yêu cầu \texttt{?id=1} nhưng vẫn chặn đủ 15/15 biến thể ở Bảng~\ref{tab:bypass} và cả bốn kịch bản ở Chương~\ref{chap:demo}. Tuy vậy, PL3 vẫn chặn nhầm sổ lưu bút và PL4 còn bốn trường hợp; mỗi trường hợp cần một luật loại trừ riêng, viết đủ hẹp để không mở lỗ hổng mới. Quy trình hợp lý là chạy DetectionOnly với lưu lượng thật, đọc audit log tìm luật khớp với yêu cầu hợp lệ, viết luật loại trừ hẹp nhất có thể, rồi chạy lại bộ tấn công. Việc này phải làm lại cho từng ứng dụng và sau mỗi lần nâng cấp CRS; CRS chỉ có sẵn gói loại trừ cho một số ứng dụng phổ biến như WordPress, Drupal, Nextcloud.

\subsection{Chi phí hiệu năng}
Bảng~\ref{tab:bench} là kết quả đo bằng ApacheBench với 10 kết nối đồng thời, mỗi cấu hình 3000 yêu cầu, chạy 5 vòng xen kẽ và lấy trung vị. Yêu cầu đo là một tệp CSS tĩnh kèm ba tham số: WAF vẫn kiểm tra đầy đủ, còn thời gian xử lý của backend gần bằng 0, nên chênh lệch phản ánh chi phí của WAF.

\begin{table}[!htb]
\centering
\caption{Hiệu năng theo cấu hình (trung vị 5 vòng, \texttt{ab -n 3000 -c 10})}
\label{tab:bench}
\footnotesize
\begin{tabular}{@{}lrrr@{}}
\toprule
\textbf{Cấu hình} & \textbf{Yêu cầu/giây} & \textbf{Thời gian TB (ms)} & \textbf{Thêm so với \texttt{Off} (ms)} \\ \midrule
DVWA trực tiếp & 14\,555 & 0,69 & -- \\
Nginx, ModSecurity tắt (\texttt{Off}) & 5\,944 & 1,68 & 0 \\
ModSecurity + CRS, PL1 & 2\,221 & 4,50 & 2,82 \\
ModSecurity + CRS, PL4 & 1\,667 & 6,00 & 4,32 \\
\bottomrule
\end{tabular}
\end{table}

Mỗi yêu cầu tốn thêm khoảng 2,8~ms ở PL1 và 4,3~ms ở PL4 so với chỉ dùng nginx làm proxy; thông lượng giảm khoảng 63\% và 72\%. Với một trang động mất hàng chục ms để sinh ra, độ trễ thêm này là nhỏ, nhưng dưới tải cao WAF tiêu tốn CPU đáng kể và phải được tính khi định cỡ hệ thống. Số liệu đo trên một máy 8 nhân chạy chung các dịch vụ khác, chỉ có ý nghĩa so sánh tương đối.

\section{Hướng phát triển}
\begin{itemize}
  \item \textbf{Tập trung nhật ký:} đẩy audit log JSON về Wazuh hoặc Loki/Grafana của nhóm khác, giữ tối thiểu 03 tháng và cảnh báo theo luật 949110, để đáp ứng mục 7.2.1.4 của TCVN 11930:2017.
  \item \textbf{Chống dò quét và giới hạn tốc độ:} kết hợp \texttt{limit\_req}/\texttt{limit\_conn} của nginx hoặc CrowdSec để chặn theo IP nguồn, bù cho phần CRS không làm được (giới hạn kết nối, brute force).
  \item \textbf{Thử Coraza:} chạy cùng bộ CRS và cùng \texttt{attack.sh} trên Coraza (Caddy hoặc Envoy) để so sánh tỷ lệ chặn và hiệu năng với ModSecurity v3.
  \item \textbf{Mở rộng kịch bản:} RFI, CSRF, SSRF, XXE và DVWA ở mức \textit{medium}/\textit{high}; đo hiệu năng trên máy riêng với trang động thay vì tệp tĩnh.
\end{itemize}

\section{Kết luận}
Nhóm đã dựng được lab cô lập Nginx + ModSecurity v3 + OWASP CRS 4.30.0 bảo vệ DVWA, khởi chạy lại bằng một lệnh \texttt{docker compose up -d} với image ghim theo digest, và tự động hóa việc tấn công, so sánh, thu bằng chứng bằng \texttt{attack.sh}. Kết quả chính:
\begin{itemize}
  \item Ở PL1, WAF chặn cả 4/4 kịch bản SQLi, XSS, LFI, Command Injection vốn khai thác thành công khi không có WAF; ở DetectionOnly cùng các luật đó được ghi log nhưng yêu cầu vẫn đi qua.
  \item Với 15 biến thể payload, PL1 chặn 11/15. Bốn biến thể command injection dùng lệnh ngắn (\texttt{id}, \texttt{pwd}) lọt qua và chạy được; phải lên PL3 mới chặn đủ 15/15.
  \item Trên 10 thao tác hợp lệ, PL1 không chặn nhầm; PL2 chặn nhầm 1, PL3 chặn nhầm 3 (còn 1 sau một dòng luật loại trừ), PL4 chặn nhầm 6.
  \item Mỗi yêu cầu tốn thêm khoảng 2,8~ms ở PL1 và 4,3~ms ở PL4 so với nginx không bật WAF.
  \item Đối chiếu 11 điều khoản TCVN 11930:2017: đáp ứng 2, một phần 7, không đáp ứng 2. WAF phủ tốt yêu cầu chống tấn công ứng dụng nhưng cần thêm SIEM và giới hạn kết nối mới đủ cho cấp độ 3.
\end{itemize}

Từ quá trình làm lab, nhóm rút ra bốn bài học vận hành. Thứ nhất, luôn chạy DetectionOnly trước rồi mới bật chặn, vì chỉ lưu lượng thật mới cho thấy luật nào chặn nhầm. Thứ hai, xử lý chặn nhầm bằng luật loại trừ hẹp nhất có thể (một luật, một biến), không tắt cả nhóm luật. Thứ ba, chọn paranoia level là chọn giữa bỏ sót và chặn nhầm: PL1 phù hợp đa số ứng dụng, chỉ nên tăng khi có người duy trì luật loại trừ. Cuối cùng, WAF là một lớp trong phòng thủ nhiều lớp, giúp giảm rủi ro trong lúc chờ sửa mã chứ không thay thế việc viết ứng dụng an toàn.
